\subsection{Equivariant functions}

In this subsection, we fix a closed subgroup H of G. We denote by \(\varpi\) the canonical projection of G onto G/H.

In addition to the Haar measure \(\mu\) on G, we fix a left Haar measure \(\beta\) on H.

By INT, VII, p. 56, § 2, no. 5, th. 2, there exists a continuous function \(\kappa\colon\mathbf{G}\to\mathbf{R}_{+}^{*}\) such that \(\kappa(xh)=\Delta_{\mathrm{H}}(h)\Delta_{\mathrm{G}}(h)^{-1}\kappa(x)\) for every \((x,h)\in\mathbf{G}\times\mathbf{H}\) . We fix such a function \(\kappa\) . We denote \(\nu\) the measure \((\kappa\cdot\mu)/\beta\) on G/H; by loc. cit., it is a non-zero positive measure quasi-invariant under G. Its support is equal to G/H (INT,

VII, p. 10, § 1, no. 1). By INT, VII, p. 43, § 2, no. 2, prop. 4, the measure \(\nu\) is the unique measure on G/H such that the measure \(\nu^{\sharp}\) on G is equal to \(\kappa \cdot \mu\) . We equip the space G/H with the measure \(\nu\) (so that we write, for example, \(\mathcal{L}^{p}(\mathrm{G / H}) = \mathcal{L}^{p}(\mathrm{G / H},\nu)\) ).

 We say that a set \(S \subset G\) is negligible modulo H if \(\varpi(S)\) is \(\nu\) -negligible. This condition does not depend on the choice of Haar measures on G and on H. It implies that S is locally \(\mu\) -negligible (INT, VII, p. 47, § 2, no. 3, prop. 6, a)). We say that a property of an element of G is true almost everywhere modulo H if the set of elements of G for which this property does not hold is negligible modulo H.

Let \(\pi\) a unitary representation of H in a Hilbert space E. We denote by \(\mathcal{F}_{\pi}(\mathrm{G})\) the vector space of functions \(f\) on G with values in E such that \(f(xh) = \pi(h)f(x)\) for all \((x,h) \in \mathrm{G} \times \mathrm{H}\) . We say that the elements of \(\mathcal{F}_{\pi}(\mathrm{G})\) are the \(\pi\) -equivariant functions on G.

The space \(\mathcal{F}_{1}(\mathrm{G})\) associated with the trivial representation of G on C is identified with the space \(\mathcal{F}(\mathrm{G / H})\) of complex-valued functions on G/H by the map \(f\mapsto f\circ \varpi\) of \(\mathcal{F}(\mathrm{G / H})\) in \(\mathcal{F}_1(\mathrm{G})\) .

For any function \(f\) in \(\mathcal{F}_{\pi}(\mathrm{G})\) , the function \(\|f\|\) belongs to \(\mathcal{F}_{1}(\mathrm{G})\) since \(\pi\) is unitary. This allows us to identify \(\|f\|\) with a function on \(\mathrm{G/H}\) , and we shall write, for example, \(\|f(x\mathrm{H})\|\) for the value of this function at an element \(x\mathrm{H}\) of \(\mathrm{G/H}\) .

A function \(f\) of \(\mathcal{F}_{\pi}(\mathrm{G})\) will be said to vanish outside a compact set modulo H if \(\|f\|\) vanishes outside a compact subset of G/H. This is equivalent to saying that there exists a compact subset K of G such that \(f\) vanishes outside \(\mathrm{K} \cdot \mathrm{H}\) (TG, III, p. 33, prop. 10).

We denote \(\mathcal{K}_{\pi}(\mathrm{G})\) the space of continuous functions on G belonging to \(\mathcal{F}_{\pi}(\mathrm{G})\) which have compact support modulo H.

A space analogous to \(\mathcal{H}_{\pi}(\mathrm{G})\) was defined in INT, VII, p. 39, § 2, no. 1, when \(\pi\) is a continuous homomorphism of H into \(\mathbf{R}_{+}^{*}\) .

Let \(p \in [1, +\infty[\) . For every \(f \in \mathcal{F}_{\pi}(\mathrm{G})\) , we denote

\[
\mathrm{N} _ {p} (f) = \left(\int_ {\mathrm{G/H}} ^ {*} \| f \| ^ {p} d \nu\right) ^ {1 / p}.
\]

It is a real number or \(+\infty\) . We denote \(\mathcal{F}_{\pi}^{p}(\mathrm{G},\nu)\) , or simply \(\mathcal{F}_{\pi}^{p}(\mathrm{G})\) , the subspace of functions \(f\in \mathcal{F}_{\pi}(\mathrm{G})\) such that \(\mathrm{N}_p(f)\)

is finite. The space \(\mathcal{F}_{\pi}^{p}(G)\) equipped with the map \(N_{p}\) is a semi-normed space.

The space \(\mathcal{K}_{\pi}(\mathrm{G})\) is contained in \(\mathcal{F}_{\pi}^{p}(\mathrm{G})\) . Its closure in \(\mathcal{F}_{\pi}^{p}(\mathrm{G})\) is denoted by \(\mathcal{L}_{\pi}^{p}(\mathrm{G},\nu)\) , or simply \(\mathcal{L}_{\pi}^{p}(\mathrm{G})\) ; we say that it is the space of \(\pi\) -equivariant functions on G of power \(p^{e}\) integrable modulo H.

The assertions of the following proposition, when \(\pi\) is the trivial representation of dimension 1, are consequences of INT, IV, p. 128, § 3, no. 3, prop. 6 and p. 131, § 3, no. 4, th. 3.

PROPOSITION 3. --- a) Let \((f_n)_{n \in \mathbf{N}}\) a sequence in \(\mathcal{F}_\pi^p(\mathrm{G})\) such that the series

\[
\sum_ {n = 0} ^ {+ \infty} \mathrm{N} _ {p} (f _ {n})
\]

converges. Then the series with general term \(f_n(g)\) is absolutely convergent for \(g\) outside a set T which is negligible modulo H. Let \(f\) be the function on G with values in E which is equal to the sum of this series on G---T and equal to 0 on T. We then have \(f \in \mathcal{F}_\pi^p(\mathrm{G})\) and the series with general term \(f_n\) converges to \(f\) in the space \(\mathcal{F}_\pi^p(\mathrm{G})\) ;

\begin{enumerate}
\def\labelenumi{\alph{enumi})}
\setcounter{enumi}{1}
\item
  Let \((f_n)\) a sequence converging in \(\mathcal{L}_{\pi}^{p}(\mathrm{G})\) to a function \(f\) . There exists a sequence \((f_{n_k})\) extracted from \((f_n)\) such that \(f_{n_k}(g)\) converges to \(f(g)\) for all \(g\) outside a set negligible modulo H;
\item
  The semi-normed spaces \(\mathcal{F}_{\pi}^{p}(\mathrm{G})\) and \(\mathcal{L}_{\pi}^{p}(\mathrm{G})\) are complete.
\end{enumerate}

In the assertion \(a\) ), saying that the series with general term \(f_n\) converges to \(f\) in the space \(\mathcal{F}_{\pi}^{p}(\mathrm{G})\) means that the sequence of partial sums \(f_0 + \cdots + f_n\) converges to \(f\) in \(\mathcal{F}_{\pi}^{p}(\mathrm{G})\) . One then also says that \(f\) is a sum of this series.

Let us prove \(a\) ). By prop. 6 of INT, IV, p. 128, § 3, no. 3, there exists a set \(\nu\) -negligible S ⊂ G/H such that the series with general term \(\|f_n(gH)\|\) converges absolutely for gH \(\notin\) . S. Moreover, the function h which is equal to the sum of this series for gH \(\notin\) . S and which is zero for gH \(\in\) . S satisfies N \(_p(h)\) \textless{} + \(\infty\) .

The set T = \(\overline{\varpi}(S)\) is negligible modulo H. For every \(g \notin T\) , the series with general term \(f_{n}(g)\) is absolutely convergent in E. Let us define \(f(g) = \sum f_{n}(g)\) for \(g \notin T\) and \(f(g) = 0\) otherwise. We have \(f \in \mathcal{F}_{\pi}(G)\) . Note that \(\|f(gH)\| \leqslant h(gH)\) for all \(g \in G\) , whence \(N_{p}(f) \leqslant\)

\(N_{p}(h) < +\infty\) . We therefore have \(f \in \mathcal{F}_{\pi}^{p}(G)\) . Similarly, it follows that

\[
\mathrm{N} _ {p} \Big (f - \sum_ {n = 0} ^ {k} f _ {n} \Big) \leqslant \sum_ {n = k + 1} ^ {+ \infty} \mathrm{N} _ {p} (f _ {n})
\]

for all \(k \in \mathbf{N}\) , hence the series \(\sum f_n\) converges to \(f\) in \(\mathcal{F}_{\pi}^{p}(\mathrm{G})\) . The assertion \(a)\) is proved.

Let us prove \(b\) ). The sequence \((\|f_n - f\|)_{n \in \mathbf{N}}\) converges to 0 in the space \(\mathcal{L}^p(\mathrm{G}/\mathrm{H})\) . By INT, IV, p. 131, § 4, no. 3, th. 3, there exists a subsequence \((\|f_{n_k} - f\|)_{k \in \mathbf{N}}\) which converges \(\nu\) -almost everywhere to 0. The sequence \((f_{n_k}(g))_{k \in \mathbf{N}}\) then converges to \(f(g)\) for all \(g\) outside a set negligible modulo H.

Finally let us prove \(c\) ). Let \((f_n)_{n \in \mathbf{N}}\) a Cauchy sequence in \(\mathcal{F}_\pi^p(\mathrm{G})\) . There exists a strictly increasing sequence of integers \((n_k)_{k \in \mathbf{N}}\) such that \(\mathrm{N}_p(f_{n_{k+1}} - f_{n_k}) \leqslant 2^{-k}\) for all \(k \in \mathbf{N}\) . For every \(k \in \mathbf{N}\) set \(h_k = f_{n_{k+1}} - f_{n_k} \in \mathcal{F}_\pi^p(\mathrm{G})\) . By \(a\) ), the series with general term \(h_k\) converges in \(\mathcal{F}_\pi^p(\mathrm{G})\) ; denote it \(h\) its sum. For every \(\ell \in \mathbf{N}\) , it follows that

\[
f _ {n _ {0}} + \sum_ {k = 0} ^ {\ell} h _ {k} = f _ {n _ {\ell + 1}},
\]

hence \(f_{n_{0}} + h\) is a cluster point of the sequence \((f_{n})\) . The latter is therefore convergent (TG, II, p. 14, cor. 2 of prop. 5). Thus the space \(\mathcal{F}_{\pi}^{p}(\mathrm{G})\) is complete; since the space \(\mathcal{L}_{\pi}^{p}(\mathrm{G})\) is closed in \(\mathcal{F}_{\pi}^{p}(\mathrm{G})\) , it is also complete.

COROLLARY. --- Let \((f_{n})_{n\in\mathbf{N}}\) a Cauchy sequence in \(\mathcal{L}_{\pi}^{p}(\mathrm{G})\) and let \(f\in\mathcal{F}_{\pi}(\mathrm{G})\) such that \(f_{n}(g)\) converges to \(f(g)\) almost everywhere modulo H. Then \(f\in\mathcal{L}_{\pi}^{p}(\mathrm{G})\) and \((f_{n})_{n\in\mathbf{N}}\) converges to f in \(\mathcal{L}_{\pi}^{p}(\mathrm{G})\) .

Indeed, the function \(f\) coincides almost everywhere modulo H with the limit of the sequence \((f_n)\) in \(\mathcal{L}_{\pi}^{p}(\mathrm{G})\) .

We denote \(\mathrm{L}_{\pi}^{p}(\mathrm{G})\) the separated normed topological vector space associated with the seminormed space \(\mathcal{L}_{\pi}^{p}(\mathrm{G})\) ; it is a Banach space.

Let \(f\) a function on \(\mathrm{G}\) with values in \(\mathrm{E}\) . We denote \(\mathrm{S}_{f}\) the set of \(x\in \mathrm{G}\) such that the function \(h\mapsto f(xh)\) onto \(\mathrm{H}\) does not belong to \(\mathcal{L}_{\mathrm{E}}^{1}(\mathrm{H})\) . We have \(\mathrm{S}_f\cdot h = \mathrm{S}_f\) for all \(h\in \mathrm{H}\) .

Let \(x \in \mathrm{G} - \mathrm{S}_f\) . Since \(\pi\) is a unitary representation, the function \(h \mapsto \pi(h)^* f(xh)\) is measurable (as a composition of the map \(h \mapsto (h, f(xh))\) from H to \(\mathrm{H} \times \mathrm{E}\) which is then measurable and of the continuous map \((g,x)\mapsto\pi(g)^{*}x\) from H × E to E, cf. INT, IV, p. 174, § 5, no. 3, th. 1), and it is integrable on H.

One defines a function \(f^{\pi}\) on G by setting

\[
f ^ {\pi} (x) = \int_ {\mathrm{H}} \pi (h) ^ {*} f (x h) d \beta (h)\tag{3}
\]

for \(x \in \mathrm{G} - \mathrm{S}_f\) and \(f^\pi(x) = 0\) if \(x \in \mathrm{S}_f\) .

PROPOSITION 4. --- Let \(f \in \mathcal{L}_{\mathrm{E}}^{1}(\mathrm{G})\) .

\begin{enumerate}
\def\labelenumi{\alph{enumi})}
\item
  The set \(\mathrm{S}_f\) is negligible modulo H and \(f^\pi \in \mathcal{F}_\pi(\mathrm{G})\) ;
\item
  Let C be a compact subset of G. If f is continuous and supported in C, then \(S_{f}\) is empty, the function \(f^{\pi}\) belongs to \(\mathcal{H}_{\pi}(\mathrm{G})\) and its support is contained in \(C \cdot H\) .
\end{enumerate}

The first part of \(a\) ) follows from INT, VII, p. 57, § 2, no. 5, c). Let \(w \in H\) . If \(x \in G - S_f\) , then \(xw \in G - S_f\) and

\[
\begin{array}{c} f ^ {\pi} (x w) = \int_ {\mathrm{H}} \pi (h) ^ {*} f (x w h) d \beta (h) \\ = \int_ {\mathrm{H}} \pi (w ^ {- 1} y) ^ {*} f (x y) d \beta (y) = \pi (w) f ^ {\pi} (x), \end{array}
\]

whereas if \(x \in S_{f}\) , then \(f^{\pi}(xw) = 0 = \pi(w)^{*}f^{\pi}(x)\) . The function \(f^{\pi}\) therefore belongs to \(\mathcal{F}_{\pi}(\mathrm{G})\) .

Suppose now that \(f\) is continuous and supported in C. For every \(x \in G\) , the map \(h \mapsto f(xh)\) then belongs to \(\mathcal{K}(H)\) , hence is integrable, which proves that \(S_f = \varnothing\) .

Let us prove that \(f^{\pi}\) is continuous. Let \(x \in G\) . Let U be a relatively compact open neighborhood of \(x\) in G.

For every \(y \in \mathrm{U}\) , we have

\[
\begin{array}{l} \| f ^ {\pi} (y) - f ^ {\pi} (x) \| \leqslant \int_ {\mathrm{H}} \| f (y h) - f (x h) \|   d \beta (h) \\ \qquad = \int_ {\mathrm{H} \cap (y ^ {- 1} \mathrm{C} \cup x ^ {- 1} \mathrm{C})} \| f (y h) - f (x h) \|   d \beta (h) \\ \qquad \leqslant \beta (\mathrm{H} \cap \mathrm{U} ^ {- 1} \mathrm{C}) \sup _ {h \in \mathrm{U} ^ {- 1} \mathrm{C}} \| f (y h) - f (x h) \|, \end{array}
\]

and the continuity of \(f^{\pi}\) then follows from the uniform continuity of f on G.

Finally, if \(x \in G\) satisfies \(f^{\pi}(x) \neq 0\) , there exist \(h \in H\) such that \(f(xh) \neq 0\) thus xh belongs to C and x belongs to C·H. Since C·H is closed in G, we conclude that the support of \(f^{\pi}\) is contained in C·H.

 Let C be a compact subset of G. Let \(u \in \mathcal{K}_{+}(G)\) a function such that \(u(x) > 0\) for all \(x \in C\) . Let v be the function on G defined by

\[
v (x) = \int_ {\mathrm{H}} u (x h) d \beta (h)\tag{4}
\]

for all \(x \in G\) ; with the previous notation, we have \(v = u^{1}\) , corresponding to the trivial 1-dimensional representation of H. The function v is continuous and positive; it belongs to \(\mathcal{F}_{1}(G)\) , its support is contained in \(\operatorname{Supp}(u) \cdot H\) and we have

\[
\inf _ {x \in \mathrm{C} \cdot \mathrm{H}} v (x) > 0\tag{5}
\]

(INT, VII, p. 39--40, § 2, no. 1, prop. 1 and lemma 1, a)).

Lemma 5. --- Let \(f \in \mathcal{F}_{\pi}(\mathrm{G})\) a function \(\mu\) -measurable and zero outside of \(\mathrm{C} \cdot \mathrm{H}\) such that the function \(\|f\|\) is \(\nu\) -integrable on \(\mathrm{G}/\mathrm{H}\) . Let s be the function on \(\mathrm{G}\) with values in E such that

\[
 s (x) = \left\{ \begin{array}{l l} v (x) ^ {- 1} f (x) & \text{if } x \in \mathrm{C} \cdot \mathrm{H} \\ 0 & \text{if } x \in \mathrm{G} - \mathrm{C} \cdot \mathrm{H}. \end{array} \right.
\]

\begin{enumerate}
\def\labelenumi{\alph{enumi})}
\item
  The function s is \(\mu\) -measurable; it belongs to \(\mathcal{F}_{\pi}(\mathrm{G})\) and is zero outside of C·H;
\item
  The function us belongs to \(\mathcal{L}^1 (\mathrm{G})\) and \((us)^{\pi} = f\) almost everywhere modulo H.
\end{enumerate}

The function \(s\) vanishes outside C \(\cdot\) . H by definition; it is \(\mu\) -measurable since the function \(f\) is \(\mu\) -measurable and \(v(x) > 0\) for all \(x \in \mathrm{C} \cdot \mathrm{H}\) (INT, IV, p. 193, § 5, no. 10, prop. 16). We have \(s \in \mathcal{F}_{\pi}(\mathrm{G})\) because \(v \in \mathcal{F}_1(\mathrm{G})\) .

The function \(f\) is locally \(\mu\) -integrable, since \(\|f\|\) is a function \(\nu\) -integrable (INT, VII, p. 47, § 2, no. 3, prop. 6, c), noting that the measure \(\nu^{\sharp} = \kappa \cdot \mu\) is equivalent to \(\mu\) ), hence the function \(s\) is likewise by formula (5). The function \(us\) is measurable and zero outside a compact subset of G; since \(u\) is bounded, it follows that \(us\) is \(\mu\) -integrable.

For every \(x\in \mathrm{G - S}_{us}\) , we have

\[
\begin{array}{c} (u s) ^ {\pi} (x) = \int_ {\mathrm{H}} \pi (h) ^ {*} (u (x h) s (x h)) d \beta (h) \\ = \int_ {\mathrm{H}} u (x h) s (x) d \beta (h) = v (x) s (x), \end{array}
\]

since \(s(xh) = \pi(h)s(x)\) ; the last assertion follows from prop. 4, a).

When \(H = \{e\}\) and \(\pi\) is of dimension 1, the following proposition is none other than INT, IV, p. 184, § 5, no. 6, th. 5.

PROPOSITION 5. --- Let \(p \in [1, +\infty[\) . The space \(\mathcal{L}_{\pi}^{p}(\mathrm{G})\) is the space of functions \(f \in \mathcal{F}_{\pi}(\mathrm{G})\) such that \(f\) is \(\mu\) -measurable and that the function \(\|f\|\) belongs to \(\mathcal{L}^{p}(\mathrm{G}/\mathrm{H})\) .

Let \(f \in \mathcal{L}_{\pi}^{p}(G)\) . It is the limit in \(\mathcal{L}_{\pi}^{p}(G)\) of a sequence of elements of \(\mathcal{K}_{\pi}(G)\) . By prop. 3, b) and Egoroff's theorem (INT, IV, p. 175, § 5, no. 4, th. 2), the function \(f\) is therefore \(\mu\) -measurable; consequently, the function \(\|f\|\) on G/H is \(\nu\) -measurable (INT, VII, p. 47, § 2, no. 3, prop. 6, b)). Since \(N_{p}(f)\) is finite, the function \(\|f\|\) belongs to \(\mathcal{L}^{p}(G/H)\) (INT, IV, p. 184, § 5, no. 6, th. 5).

 Let us prove the converse assertion. Let \(f\) a function \(\mu\) -measurable belonging to \(\mathcal{F}_{\pi}(\mathrm{G})\) such that \(\|f\|\in\mathcal{L}^{p}(\mathrm{G}/\mathrm{H})\) . Let \(\varepsilon>0\) . Let us prove that there exists \(\widetilde{f}\in\mathcal{K}_{\pi}(\mathrm{G})\) such that

\[
\mathrm{N} _ {p} (f - \widetilde {f}) ^ {p} = \int_ {\mathrm{G/H}} ^ {*} \| f - \widetilde {f} \| ^ {p} d \nu \leqslant \varepsilon ,
\]

which will conclude the proof.

 Suppose first that there exists a compact subset C of G such that \(f\) vanishes outside \(\mathrm{C} \cdot \mathrm{H}\) . Let \(u \in \mathcal{K}_{+}(\mathrm{G})\) a function such that \(u(x) > 0\) for all \(x \in \mathrm{C}\) and let us define \(v = u^{1}\) as above. Denote \(\varphi\) the characteristic function of the support of \(u\) .

Let q be the conjugate exponent of p. If p = 1, let w be the constant function on G equal to \(\sup_{x \in G} u(x)\) . If p \textgreater{} 1, we define

\[
w (x) = \left(\int_ {\mathrm{H}} u (x h) ^ {q} d \beta (h)\right) ^ {1 / q}
\]

 for all \(x \in G\) . In all cases, the function w is continuous and positive; it belongs to \(\mathcal{K}_{1}(G)\) (Prop. 4 applied to the trivial representation of dimension 1 and to the function \(u^{q}\) ) hence it is bounded on G. We set \(W = \sup_{x \in G} w(x)\) .

 Let \(s\) the function on \(\mathbf{G}\) with values in \(\mathbf{E}\) defined by Lemma 5 applied to \(f\) . Set \(g = \varphi s\) . The function \(g\) is \(\mu\) -measurable, and satisfies \(\| g \| \leqslant \| f \| / \inf_{x \in \mathrm{C} \cdot \mathrm{H}} v(x)\) . Since \(g\) is zero outside the support of \(u\) and that \(\kappa\) is continuous, we have \(g \in \mathcal{L}_{\mathrm{E}}^{p}(\mathrm{G}, \kappa \cdot \mu)\) . Let \(\widetilde{g}\) a function

in \(\mathcal{K}_{\mathrm{E}}(\mathrm{G})\) such that

\[
\int_ {\mathrm{G}} ^ {*} \| g - \widetilde {g} \| ^ {p} \kappa d \mu \leqslant \frac {\varepsilon}{\mathrm{W} ^ {p}}.
\]

We have us = ug, hence \(f = (us)^{\pi} = (ug)^{\pi}\) almost everywhere modulo H (lemma 5, b)). Let \(\widetilde{f} = (u\widetilde{g})^{\pi}\) ; we have \(\widetilde{f} \in \mathcal{K}_{\pi}(\mathrm{G})\) (prop. 4, b)). For every \(x \in G - S_{ug}\) , we obtain

\[
\| (u g) ^ {\pi} (x) - (u \widetilde {g}) ^ {\pi} (x) \| \leqslant \int_ {\mathrm{H}} ^ {*} u (x h) \| g (x h) - \widetilde {g} (x h) \| d \beta (h)
\]

whence

\[
\| (u g) ^ {\pi} (x) - (u \widetilde {g}) ^ {\pi} (x) \| ^ {p} \leqslant w (x) ^ {p} \int_ {\mathrm{H}} ^ {*} \| g (x h) - \widetilde {g} (x h) \| ^ {p} d \beta (h),
\]

by Hölder's inequality in the case \( p \) \textgreater{} 1. Since \(S_{ug}\) is negligible modulo H, it then follows that

\[
\begin{array}{c} \int_ {\mathrm{G/H}} ^ {*} \| f - \widetilde {f} \| ^ {p} d \nu \leqslant \mathrm{W} ^ {p} \int_ {\mathrm{G/H}} ^ {*} \Big (\int_ {\mathrm{H}} ^ {*} \| g (x h) - \widetilde {g} (x h) \| d \beta (h) \Big) ^ {p} d \nu (x \mathrm{H}) \\ = \mathrm{W} ^ {p} \int_ {\mathrm{G}} ^ {*} \| g - \widetilde {g} \| ^ {p} \kappa d \mu \leqslant \varepsilon \end{array}
\]

according to INT, VII, p. 46, § 2, no. 3, prop. 5, b), which is applicable because \(g - \widetilde{g}\) vanishes outside a compact subset of G. This implies the required property when \(f\) is zero outside a compact subset modulo H.

Now consider the general case. Since \(\|f\|\in\mathcal{L}^{p}(G/H)\) by hypothesis, there exists a compact subset L of G/H such that

\[
\int_ {(\mathrm{G} / \mathrm{H}) - \mathrm{L}} ^ {*} \| f \| ^ {p} d \nu \leqslant \frac {\varepsilon}{2}
\]

(cf. INT, IV, p. 152, § 4, no. 6, th. 4). Let \(\varphi_{\mathrm{L}}\) the characteristic function of L and set \(f_{\mathrm{L}} = (\varphi_{\mathrm{L}} \circ \varpi)f\) . We have

\[
\int_ {\mathrm{G} / \mathrm{H}} ^ {*} \| f - f _ {\mathrm{L}} \| ^ {p} d \nu = \int_ {(\mathrm{G} / \mathrm{H}) - \mathrm{L}} ^ {*} \| f \| ^ {p} d \nu \leqslant \frac {\varepsilon}{2}.
\]

The function \(f_{\mathrm{L}}\) is \(\mu\) -measurable and zero outside a compact set modulo H. It belongs to \(\mathcal{L}_{\pi}^{p}(\mathrm{G})\) therefore, by the preceding case, there exists \(\widetilde{f} \in \mathcal{K}_{\pi}(\mathrm{G})\) such that \(\mathrm{N}_p(f_{\mathrm{L}} - \widetilde{f}) \leqslant (\frac{\varepsilon}{2})^{1/p}\) , whence

\[
\int_ {\mathrm{G/H}} ^ {*} \| f - \widetilde {f} \| ^ {p} d \nu \leqslant \varepsilon ,
\]

as desired.

 Let us consider the case \(p = 2\) . For \(f_1\) and \(f_2\) in \(\mathcal{F}_{\pi}(\mathrm{G})\) , the function \(x \mapsto \langle f_1(x) | f_2(x) \rangle\) belongs to \(\mathcal{F}_1(\mathrm{G})\) since the representation \(\pi\) is unitary, and therefore defines by passing to the quotient a function on \(\mathrm{G / H}\) which we identify as before with \(\langle f_1 | f_2 \rangle\) . We have the bound \(|\langle f_1 | f_2 \rangle| \leqslant \| f_1 \| \| f_2 \|\) in \(\mathcal{F}(\mathrm{G / H})\) .

If \(f_1\) and \(f_2\) belong to \(\mathcal{L}_\pi^2(\mathrm{G})\) , then the function \(\langle f_1 | f_2 \rangle\) belongs to \(\mathcal{L}_1^1(\mathrm{G})\) . In particular, it is integrable on \(\mathrm{G/H}\) . The map

\[
(f _ {1}, f _ {2}) \mapsto \int_ {\mathrm{G} / \mathrm{H}} \langle f _ {1} | f _ {2} \rangle d \nu
\]

is a positive hermitian form on \(\mathcal{L}_{\pi}^{2}(\mathrm{G})\) ; the associated semi-norm is the semi-norm \(N_{2}\) . In particular, the space \(\mathcal{L}_{\pi}^{2}(\mathrm{G})\) is a pre-Hilbert space and \(\mathrm{L}_{\pi}^{2}(\mathrm{G})\) is the Hilbert space associated with \(\mathcal{L}_{\pi}^{2}(\mathrm{G})\) .

